\section{Labor heterogeneity and reallocation}
\label{sec:rewrite-labor-heterogeneity}

Does the wage-growth result require all workers to have the same relation
to AI? I consider two activities, one in which labor and AI are complements
and another in which they are substitutes. Both produce the same final good,
so their outputs are perfect substitutes. This assumption allows aggregate
production to expand without requiring both activities to expand together.
I characterize long-run outcomes conditional on an equilibrium with the
stated limiting properties, rather than establish a new existence result.

For $j\in\{S,C\}$, production is
\begin{align}
 Y&=Y_S+Y_C,\qquad Y_j=K_j^\alpha Z_j^{1-\alpha},
 \label{eq:rewrite-heterogeneous-output}\\
 Z_j&=\left[\omega_{Lj}(AL_j)^{(\sigma_j-1)/\sigma_j}
 +\omega_{Xj}X_j^{(\sigma_j-1)/\sigma_j}\right]^{\sigma_j/(\sigma_j-1)},
 \qquad 0<\sigma_C<1<\sigma_S.
 \label{eq:rewrite-heterogeneous-ces}
\end{align}
The CES weights are positive and sum to one within each activity.
Capital is freely mobile, $K_S+K_C=K$, and both activities use the same
AI efficiency, $X_j=BU_j$. The developer charges a common price $p_X$.
The representative household pools the two labor incomes and retains the
preferences and asset ownership of Section~\ref{sec:rewrite-model}.
Its wage income is now $w_SL_S+w_CL_C$, and the resource constraint is
$C+\dot K+\delta K+(X_S+X_C)/B+M=Y$.
The research technology is unchanged.

Optimal capital allocation gives
\begin{equation}
 \frac{K_j}{K}=\frac{Z_j}{Z_S+Z_C},\qquad
 Y=K^\alpha(Z_S+Z_C)^{1-\alpha}.
 \label{eq:rewrite-heterogeneous-capital}
\end{equation}
Writing $Q=(1-\alpha)[K/(Z_S+Z_C)]^\alpha$ for the marginal value of
either composite, the input conditions in an active activity are
\begin{equation}
 p_X=Q\frac{\partial Z_j}{\partial X_j},\qquad
 w_j=AQ\omega_{Lj}\left(\frac{Z_j}{AL_j}\right)^{1/\sigma_j}.
 \label{eq:rewrite-heterogeneous-prices}
\end{equation}
The developer faces demand after capital and AI services have been allocated
between activities. It takes aggregate capital and the exogenous paths as
given, as before, and maximizes operating profit globally. The original
single-activity condition~\eqref{eq:rewrite-monopoly-foc} cannot be used
with either sector's demand elasticity in place of the joint elasticity.
Appendix~\ref{app:rewrite-labor-heterogeneity} derives the joint demand
and the limiting global supply choice.

\subsection{Sector-specific labor}
\label{subsec:rewrite-sector-specific-labor}

First, let $L_j=\ell_jN$, where the positive shares $\ell_S$ and $\ell_C$
are fixed and sum to one. Workers cannot change activities, although capital
can. Let $\mathcal R_S(B)$ denote the limiting monopoly return generated
by activity $S$:
\begin{equation}
 \mathcal R_S(B)=\alpha\left[(1-\alpha)^2B
 \omega_{XS}^{\sigma_S/(\sigma_S-1)}\right]^{(1-\alpha)/\alpha}-\delta.
 \label{eq:rewrite-heterogeneous-return}
\end{equation}
This is the return in equation~\eqref{eq:rewrite-ai-terminal-ratio}
with the substitution parameters of activity $S$.
In the following results, growth denotes the long-run average
$\lim_{t\to\infty}t^{-1}\log(V_t/V_0)$ for the indicated positive variable.
Instantaneous growth has the same limit whenever that limit exists.

\begin{proposition}[Wage growth with sector-specific labor]
\label{prop:rewrite-heterogeneous-wages}
Consider a monopoly equilibrium with $K/(AN)\to\infty$,
$B\to B_\infty\in(0,\overline B]$, and $C/K$ converging to a positive
constant. If $\mathcal R_S(B_\infty)>\rho+\gamma$, output per worker
has long-run growth $g=\mathcal R_S(B_\infty)-\rho>\gamma$, and
\begin{equation}
 g_{w_S}=\left(1-\frac1{\sigma_S}\right)\gamma+\frac g{\sigma_S},
 \qquad g_{w_C}=\gamma.
 \label{eq:rewrite-heterogeneous-wage-growth}
\end{equation}
The wage ratio $w_S/w_C$ diverges. Both labor-income shares in aggregate
output converge to zero, as does $Y_C/Y$, although activity $C$ remains active.
\end{proposition}
\begin{proof}
See Appendix~\ref{app:rewrite-labor-heterogeneity}.
\end{proof}

Substitution allows capital and AI services in $S$ to expand faster than its
effective labor endowment, sustaining wage growth above $\gamma$.
Complementarity keeps $C$ tied to the scale of its effective labor, so its
wages grow at $\gamma$. Thus workers in the activity where AI substitutes
for labor experience faster wage growth. Neither group keeps pace with
aggregate output per worker. The declining output share of $C$ does not
mean that its output or wages decline in levels, and these growth comparisons
do not establish welfare gains or the sign of wage changes during the transition.

The characterization allows $B_\infty<\overline B$ rather than importing
the efficiency-convergence result from the original model.
Under competitive provision with constant $B=B_0$, the same wage-growth
expressions apply to an equilibrium with these limiting properties, using
the marginal-cost return obtained by replacing $(1-\alpha)^2$ in
equation~\eqref{eq:rewrite-heterogeneous-return} with $1-\alpha$.

\subsection{Mobile labor}
\label{subsec:rewrite-mobile-labor}

Now let workers be identical and freely mobile, with $L_S+L_C=N$.
If both activities employ labor, their wages must coincide. The divergent
wage ratio in Proposition~\ref{prop:rewrite-heterogeneous-wages} therefore
cannot persist with fixed positive employment shares in both activities.

\Needspace{13\baselineskip}
\begin{corollary}[Reallocation to the substituting activity]
\label{cor:rewrite-mobile-labor}
Consider a monopoly equilibrium with mobile labor, $K/(AN)\to\infty$,
$B\to B_\infty\in(0,\overline B]$, and $C/K$ converging to a positive
constant. If $\mathcal R_S(B_\infty)>\rho+\gamma$, then after a finite
date all workers and productive inputs are allocated to $S$.
Long-run output-per-worker and wage growth are
\begin{equation}
 g=\mathcal R_S(B_\infty)-\rho>\gamma,\qquad
 g_w=\left(1-\frac1{\sigma_S}\right)\gamma+\frac g{\sigma_S}.
 \label{eq:rewrite-mobile-wage-growth}
\end{equation}
\end{corollary}
\begin{proof}
See Appendix~\ref{app:rewrite-labor-heterogeneity}.
\end{proof}

With free reallocation, accelerated wage growth is not restricted to workers
initially employed in $S$. They all eventually work there, and the
single-activity growth and wage formulas are recovered. This conclusion
uses both labor mobility and perfect substitution between the products.
If the product of $C$ were essential to final production, the economy could
not abandon that activity in this way.
